\section{System Architecture Details}
\label{app:architecture}

\subsection{Complete Technology Stack}
\begin{itemize}
    \item \textbf{Frontend:} React 18.2.0, TypeScript 5.3.3, Vite 5.0.0, TailwindCSS 3.4.0
    \item \textbf{Backend:} Python 3.12, FastAPI 0.109.0, PyTorch 2.6.0, CUDA 12.4
    \item \textbf{AI Models:} YOLOv11-nano (Ultralytics 8.3.0), MediaPipe 0.10.9, VideoMAE (Hugging Face Transformers 4.45.0)
    \item \textbf{Training:} Hugging Face Accelerate 1.11.0, scikit-learn 1.5.0
    \item \textbf{Hardware:} NVIDIA GeForce RTX 4080 SUPER (16GB VRAM) for training
\end{itemize}

\subsection{System Pipeline}
The complete processing pipeline operates as follows:
\begin{enumerate}
    \item Video upload via React frontend (drag-and-drop interface)
    \item Backend receives video file and stores temporarily
    \item YOLOv11 detects player(s) in each frame, extracts largest bounding box
    \item MediaPipe extracts 33 body keypoints from player region
    \item VideoMAE processes sequence of 16 frames for action classification
    \item Performance metrics engine calculates jump height, movement speed, form score, reaction time, pose stability
    \item Results returned as JSON to frontend for visualization
\end{enumerate}

\section{Dataset Statistics}
\label{app:dataset}

\subsection{Video Distribution by Category}
\begin{table}[h]
\centering
\caption{Dataset composition by action category}
\begin{tabular}{lcc}
\toprule
Action Category & Number of Videos & Percentage \\
\midrule
Free Throw Shot & 23 & 34.3\% \\
3-Point Shot & 24 & 35.8\% \\
Dribbling & 11 & 16.4\% \\
2-Point Shot & 6 & 9.0\% \\
Idle & 2 & 3.0\% \\
Defense & 1 & 1.5\% \\
Passing & 0 & 0.0\% \\
\midrule
\textbf{Total} & \textbf{67} & \textbf{100\%} \\
\bottomrule
\end{tabular}
\label{tab:dataset}
\end{table}

\subsection{Dataset Splits}
Training: 46 videos (68.7\%), Validation: 10 videos (14.9\%), Test: 11 videos (16.4\%)

\subsection{Video Characteristics}
\begin{itemize}
    \item Average duration: 5--10 seconds per clip
    \item Resolution: Variable (mobile phone recordings, typically 720p--1080p)
    \item Frame rate: 24--30 FPS
    \item Recording environment: Indoor and outdoor courts, natural lighting
    \item Equipment: Standard mobile phones (iPhone, Android devices)
\end{itemize}

\section{Model Training Details}
\label{app:training}

\subsection{Training Hyperparameters}
\begin{itemize}
    \item Base model: \texttt{MCG-NJU/videomae-base-finetuned-kinetics}
    \item Number of epochs: 192 (with early stopping)
    \item Batch size: 4 (limited by GPU memory)
    \item Learning rate: $1 \times 10^{-4}$
    \item Optimizer: AdamW with weight decay 0.01
    \item Warmup steps: 500
    \item Mixed precision: FP16 enabled
    \item Evaluation strategy: Per epoch on validation set
    \item Early stopping patience: 5 epochs
    \item Metric for best model: F1 score
\end{itemize}

\subsection{Training Results}
\begin{itemize}
    \item Final test accuracy: 90.9\%
    \item Training samples: 46 videos
    \item Validation samples: 10 videos
    \item Test samples: 11 videos
    \item Total training time: Approximately 4--6 hours on RTX 4080 SUPER
    \item Model size: 340MB (PyTorch state dict)
\end{itemize}

\section{Action Categories}
\label{app:actions}

The system classifies basketball actions into 7 categories:

\begin{enumerate}
    \item \textbf{Free Throw Shot:} Stationary shot from free throw line (4.6m from basket)
    \item \textbf{2-Point Shot:} Shot from inside 3-point arc (1.5m--6.75m from basket)
    \item \textbf{3-Point Shot:} Shot from behind 3-point line ($\geq$6.75m from basket)
    \item \textbf{Dribbling:} Ball handling while moving, maintaining ball control
    \item \textbf{Passing:} Transferring ball to teammate (chest pass, bounce pass, overhead)
    \item \textbf{Defense:} Defensive stance, guarding opponent, lateral movement
    \item \textbf{Idle:} Standing, no active action, waiting/positioning
\end{enumerate}

\section{Performance Metrics Calculation}
\label{app:metrics}

The system computes six performance metrics:

\begin{enumerate}
    \item \textbf{Jump Height:} Calculated from vertical displacement of hip keypoint during jump phase
    \item \textbf{Movement Speed:} Average velocity of center-of-mass (calculated from hip keypoints) across frames
    \item \textbf{Shooting Form Score:} Composite metric based on elbow angle, release consistency, follow-through (0--1 scale)
    \item \textbf{Reaction Time:} Time from stimulus (e.g., ball arrival) to movement initiation
    \item \textbf{Pose Stability:} Variance in keypoint positions, indicating balance and control (0--1 scale)
    \item \textbf{Energy Efficiency:} Ratio of effective movement to total movement (optional metric)
\end{enumerate}

\section{System Availability}
\label{app:availability}

The complete system, including source code, documentation, trained models, and dataset collection guidelines, is available as an open-source project. All components are licensed under permissive open-source licenses, enabling free use, modification, and distribution for academic and commercial purposes.

\section{PRISMA Flow Diagram}
\label{app:prisma}

\textbf{Note:} The PRISMA 2020 flow diagram should be included here as a figure showing:
\begin{itemize}
    \item Records identified: 102 (Scopus: 58, Web of Science: 44)
    \item Duplicates removed: 18
    \item Records screened: 84
    \item Records excluded: 50 (title/abstract screening)
    \item Full-text assessed: 34
    \item Full-text excluded: 14
    \item Studies included: 20
\end{itemize}

